\section*{NeurIPS Paper Checklist}

\begin{enumerate}

\item {\bf Claims}
    \item[] Question: Do the main claims made in the abstract and introduction accurately reflect the paper's contributions and scope?
    \item[] Answer: \answerYes{}
    \item[] Justification: The abstract claims the find/act boundary (supported by \S\ref{sec:findact}, Table~\ref{tab:findact}), the registered tier-vs-keep/drop null and the compressed-matches-full cost result at $N{=}70$ (\S\ref{sec:resolve}, Tables~\ref{tab:resolve}--\ref{tab:mcnemar}), and the $T{=}0$ run-to-run noise measurement (\S\ref{sec:resolve}). Each claim is stated with its $n$ and CI; the marginal keep/drop-vs-full comparison ($p{=}0.07$) is explicitly not claimed; scope limits (single-shot regime, oracle ceilings, one language, one model family) are stated in \S\ref{sec:threats}.

\item {\bf Limitations}
    \item[] Question: Does the paper discuss the limitations of the work performed by the authors?
    \item[] Answer: \answerYes{}
    \item[] Justification: \S\ref{sec:threats} is organized as threats-with-directions: the single-shot no-execution regime (and which way it biases each comparison), the oracle arms' asymmetric construction (why the $p{=}0.07$ comparison is not claimed), probe scope (one language, 15 classes), residual format asymmetry against the tier arm, measured non-determinism, token-unit caveats, and coverage gaps. Superseded own-code measurements are quarantined in Appendix~\ref{app:explore} and carry no claims.

\item {\bf Theory assumptions and proofs}
    \item[] Question: For each theoretical result, does the paper provide the full set of assumptions and a complete (and correct) proof?
    \item[] Answer: \answerNA{}
    \item[] Justification: The paper is empirical and contains no theoretical results.

\item {\bf Experimental result reproducibility}
    \item[] Question: Does the paper fully disclose all the information needed to reproduce the main experimental results of the paper to the extent that it affects the main claims and/or conclusions of the paper (regardless of whether the code and data are provided or not)?
    \item[] Answer: \answerYes{}
    \item[] Justification: Table~\ref{tab:setup} specifies dataset, instance criterion, models, decoding, patch protocol, harness version, validation gates, and statistics; \S\ref{sec:findact} specifies the probe protocol end to end; Appendix~\ref{app:prompt} gives the verbatim agent prompt and decoding; Appendix~\ref{app:prereg} records each pre-registration and its outcome; Appendix~\ref{app:repro} states known reproducibility limits (model alias without dated snapshot, single-draw probe cells, tokenizer proxy) and the measured $T{=}0$ run-to-run noise that bounds replication expectations.

\item {\bf Open access to data and code}
    \item[] Question: Does the paper provide open access to the data and code, with sufficient instructions to faithfully reproduce the main experimental results, as described in supplemental material?
    \item[] Answer: \answerYes{}
    \item[] Justification: Appendix~\ref{app:repro}: experiment code, frozen instance lists, committed pre-registrations, per-instance harness reports, and the figure generator (every plotted number loaded from raw result files) are released; SWE-bench Verified is public. Repository: \url{https://github.com/integrallis/act-context} (MIT license).

\item {\bf Experimental setting/details}
    \item[] Question: Does the paper specify all the training and test details (e.g., data splits, hyperparameters, how they were chosen, type of optimizer) necessary to understand the results?
    \item[] Answer: \answerYes{}
    \item[] Justification: No training occurs. All inference settings are specified: agent model and temperature, token caps, retry policy, arm construction (\S\ref{sec:rig}), instance selection rule and exclusions (Table~\ref{tab:setup}, Figure~\ref{fig:selection}), probe minting/answering/judging models and budgets (\S\ref{sec:findact}), and the registered hypotheses and decision rules (Appendix~\ref{app:prereg}).

\item {\bf Experiment statistical significance}
    \item[] Question: Does the paper report error bars suitably and correctly defined or other appropriate information about the statistical significance of the experiments?
    \item[] Answer: \answerYes{}
    \item[] Justification: Every proportion carries a Wilson 95\% CI; paired arm comparisons use exact two-sided McNemar tests on discordant pairs (Tables~\ref{tab:mcnemar}, \ref{tab:findact}); inter-judge agreement is reported as Cohen's $\kappa$; and the paper additionally \emph{measures} the run-to-run noise floor ($6/70$ outcome flips between byte-identical runs) and interprets its own small gaps against it.

\item {\bf Experiments compute resources}
    \item[] Question: For each experiment, does the paper provide sufficient information on the computer resources (type of compute workers, memory, time of execution) needed to reproduce the experiments?
    \item[] Answer: \answerYes{}
    \item[] Justification: Appendix~\ref{app:repro} (``Compute and cost''): ${\approx}210$ agent API calls (${\approx}\$20$) plus ${\approx}6$ hours of Docker evaluation on a 32-thread workstation for the registered run; ${\approx}600$ short calls (${\approx}\$5$) for the probe study; no GPUs, no training.

\item {\bf Code of ethics}
    \item[] Question: Does the research conducted in the paper conform, in every respect, with the NeurIPS Code of Ethics \url{https://neurips.cc/public/EthicsGuidelines}?
    \item[] Answer: \answerYes{}
    \item[] Justification: The work evaluates public LLM APIs on a public benchmark and public open-source repositories; no human subjects, no personal or sensitive data.

\item {\bf Broader impacts}
    \item[] Question: Does the paper discuss both potential positive societal impacts and negative societal impacts of the work performed?
    \item[] Answer: \answerYes{}
    \item[] Justification: Appendix~\ref{app:repro} (``Broader impact''): the practical implication is cheaper, lower-energy coding agents (equal resolution at a third of the context); the released assets are evaluation code and result files with no plausible misuse path.

\item {\bf Safeguards}
    \item[] Question: Does the paper describe safeguards that have been put in place for responsible release of data or models that have a high risk for misuse (e.g., pre-trained language models, image generators, or scraped datasets)?
    \item[] Answer: \answerNA{}
    \item[] Justification: No models are trained or released; the released assets are evaluation scripts, frozen instance lists, and result files over a public benchmark.

\item {\bf Licenses for existing assets}
    \item[] Question: Are the creators or original owners of assets (e.g., code, data, models), used in the paper, properly credited and are the license and terms of use explicitly mentioned and properly respected?
    \item[] Answer: \answerYes{}
    \item[] Justification: Appendix~\ref{app:repro} (``Existing assets''): SWE-bench Verified and harness \citep{swebench} (MIT); seaborn/pylint/pytest analyzed at pinned commits under their respective licenses; Anthropic, OpenAI, and Qwen (Apache-2.0) models credited and used under their terms of service; all prior work cited.

\item {\bf New assets}
    \item[] Question: Are new assets introduced in the paper well documented and is the documentation provided alongside the assets?
    \item[] Answer: \answerYes{}
    \item[] Justification: The released instrument (arm builders, expressibility auditor, patch builder, staged runners, analysis) ships with per-experiment READMEs, the committed pre-registration documents, and a figure generator that rebuilds every plot from raw result files; released under the MIT license at \url{https://github.com/integrallis/act-context}.

\item {\bf Crowdsourcing and research with human subjects}
    \item[] Question: For crowdsourcing experiments and research with human subjects, does the paper include the full text of instructions given to participants and screenshots, if applicable, as well as details about compensation (if any)?
    \item[] Answer: \answerNA{}
    \item[] Justification: No crowdsourcing and no human subjects.

\item {\bf Institutional review board (IRB) approvals or equivalent for research with human subjects}
    \item[] Question: Does the paper describe potential risks incurred by study participants, whether such risks were disclosed to the subjects, and whether Institutional Review Board (IRB) approvals (or an equivalent approval/review based on the requirements of your country or institution) were obtained?
    \item[] Answer: \answerNA{}
    \item[] Justification: No research with human subjects.

\item {\bf Declaration of LLM usage}
    \item[] Question: Does the paper describe the usage of LLMs if it is an important, original, or non-standard component of the core methods in this research?
    \item[] Answer: \answerYes{}
    \item[] Justification: Appendix~\ref{app:repro} (``LLM usage declaration''): LLMs are the evaluated subject (agent, summarizer arms, judges --- fully specified in Table~\ref{tab:setup} and \S\ref{sec:findact}); additionally, LLM-based development tooling assisted in building the experimental infrastructure and drafting, with all hypotheses, protocols, and analyses human-approved and frozen in committed pre-registrations before data collection.

\end{enumerate}
